\section{Stationary stochastic purification}\label{sec:purification55}
Write $\mathsf{Stoch}(k)$ for the compact semigroup of row-stochastic matrices of order $k$. The identity of a group inside this semigroup can be a singular idempotent; it must not be confused with the identity matrix.

\begin{lemma}[An idempotent over the physical identity]\label{lem:idempotent55}
Let $\mathcal T$ be a compact subsemigroup of $H\times\mathsf{Stoch}(k)$ whose projection onto $H$ is surjective. There is an element $e=(1,E)\in\mathcal T$ with $E^2=E$ and rank equal to the minimum matrix rank in $\mathcal T$. The monoid $e\mathcal T e$ is a compact group with identity $e$ and still projects onto $H$.
\end{lemma}
\begin{proof}
Only the integer ranks $1,\ldots,k$ occur, so their nonempty set has a minimum which is attained. Choose $(h,T)$ of that minimum rank $r$. The powers of $T$ are bounded, since they are stochastic. Thus its complex eigenvalues have modulus at most one and all Jordan blocks on the unit circle have size one: a larger peripheral Jordan block would produce polynomially growing powers. The closure of positive powers of an element of a compact group is a group: arbitrarily large positive powers approach its identity, and the preceding powers approach its inverse. Apply this observation to the compact monothetic group generated by $(h,\lambda_1,\ldots,\lambda_b)$ in $H\times(S^1)^b$, where the $\lambda_i$ are the peripheral eigenvalues of $T$. There is a net of integers tending to infinity for which $h^n\to1$ and all peripheral eigenvalues raised to the $n$th power tend to one. Along it $T^n$ tends to the spectral projection $E$ onto the peripheral subspace. This is a real stochastic idempotent. Its rank is at most $r$ and at least $r$ by minimality, so it is $r$. Nets avoid any metrizability assumption on $H$.

Every matrix $B$ in $e\mathcal T e$ satisfies $B=EBE$ and has rank $r$. Its action on the row space of $E$ is consequently invertible. Restriction embeds $e\mathcal T e$ as a compact submonoid of $H\times\mathrm{GL}(r,\R)$, with identity $(1,I_r)$; injectivity follows because the restriction determines $(xE)B=xB$ for every row $x$, by $B=EBE$, and hence determines every row of $B$. A compact submonoid of a topological group is a subgroup, by the same positive-power return argument. The inverse images of its inverses are still stochastic. Finally, sandwiching an element with physical coordinate $h$ between $e$'s leaves that coordinate unchanged. Surjectivity follows.
\end{proof}

\begin{lemma}[The recurrent simplex]\label{lem:simplex55}
If $E$ is a stochastic idempotent of rank $r$, then
\[
 \Delta_k E=\{p\in\Delta_k:pE=p\}
\]
is a simplex with exactly $r$ vertices. Every group of stochastic matrices with identity $E$ permutes these vertices.
\end{lemma}
\begin{proof}
Regard $E$ as a transition matrix. Its closed communicating classes each have a unique stationary probability row supported on that class. Every stationary row is a convex combination of these rows: mass on a transient class is zero, and stationarity within a closed irreducible class determines a one-dimensional space. For completeness, transient mass vanishes because each transient state has a positive probability of reaching a closed class in a bounded number of steps; a common bound over the finitely many states makes transient mass decrease by a fixed factor unless it is zero. The closed-class rows have disjoint supports, hence are linearly independent.

Each row of $E$ is stationary because $E^2=E$. Conversely, every stationary row equals its product with $E$. The span of the closed-class stationary rows is therefore precisely the row space of $E$, so there are $r$ classes and $r$ simplex vertices. If $B$ has a stochastic inverse relative to $E$, multiplication by $B$ and by this inverse are inverse affine maps of this simplex. Affine bijections permute its vertices.
\end{proof}

\begin{theorem}[Purification without a width or error loss]\label{thm:purification55}
For every compact-group interface in Definition~\ref{def:stationary55}, a stationary $k$-label realization of error at most $\varepsilon$ admits a permutation realization of error at most $\varepsilon$ on at most $k$ labels. There is no return-word hypothesis and no restriction on the number of connected components of $H$.
\end{theorem}
\begin{proof}
For the given machine form
\[
 \mathcal T=\overline{\{(u_w^{-1},T_w):w\in\mathcal A^*\}}
       \subset H\times\mathsf{Stoch}(k).
\]
The inverse on the physical coordinate is important: concatenation satisfies $u_{vw}^{-1}=u_v^{-1}u_w^{-1}$ and $T_{vw}=T_vT_w$. Thus $\mathcal T$ is a compact semigroup. Positive-word density makes its projection onto $H$ surjective. By continuity, its every element $(h,B)$ satisfies
\begin{equation}\label{eq:closederror55}
             \norm{\alpha_sBD_j-F_{s,j}(h^{-1})}_j\le\varepsilon
             \quad(s\in\mathcal S,\ j\in\mathcal J).
\end{equation}

Apply Lemma~\ref{lem:idempotent55} and denote the vertices of $\Delta_k E$ by $\nu_1,\ldots,\nu_r$. For each letter the element
\[
        e(u_a^{-1},T_a)e=(u_a^{-1},ET_aE)
\]
belongs to the group $e\mathcal T e$. By Lemma~\ref{lem:simplex55}, $B_a=ET_aE$ permutes the $\nu_i$. Use that permutation as the new command. If $V$ has rows $\nu_i$ and the induced row permutation is $\Pi_a$, the precise intertwining identity is $VB_a=\Pi_aV$. Write uniquely
\[
          \alpha_sE=\sum_{i=1}^r\beta_s(i)\nu_i,
          \qquad \widetilde D_j(i)=\nu_iD_j\in Y_j.
\]
The last inclusion uses convexity of the legal output set. For a nonempty word the output of the resulting $r$-label machine is $\alpha_s B_{a_1}\cdots B_{a_n}D_j$. For the empty word it is $\alpha_sED_j$. In both cases the corresponding joint element lies in $\mathcal T$ and has physical coordinate $u_w^{-1}$. Equation~\eqref{eq:closederror55} proves the error bound. Since $r\le k$, the width has not increased.
\end{proof}
The new machine need not reproduce the old machine's erroneous output word by word; it remains in the same closed error balls about the target. Nor do we assert $T_{w}=I$ on a physical return, invertibility of the original $T_a$, or deterministic initialization. These three incorrect shortcuts would invalidate the reduction.

\subsection{Finite quotients and the unrestricted stationary minimum}
For a tuple of permutation matrices $\sigma=(\Pi_a)_{a\in\mathcal A}$ of order $k$, set
\[
 \mathcal L_\sigma=
 \overline{\{(u_w^{-1},\Pi_w):w\in\mathcal A^*\}}
       \subset H\times\mathfrak S_k.
\]
The topology is the product of the compact topology on $H$ and the discrete topology on the finite group $\mathfrak S_k$. This is a compact subgroup. Its definition includes the possible discrepancy between physical relations and hidden-state relations.

\begin{theorem}[An exact stationary minimum]\label{thm:stationaryminimum55}
The value $M_\varepsilon$ is the least positive integer $k$ for which some tuple $\sigma$ of permutations of order $k$, some $\beta_s\in\Delta_k$, and some $d_j(i)\in Y_j$ satisfy
\begin{equation}\label{eq:permutationfeasibility55}
 \norm{\beta_s\Pi d_j-F_{s,j}(h^{-1})}_j\le\varepsilon
 \quad\text{for all }(h,\Pi)\in\mathcal L_\sigma\text{ and all }s,j.
\end{equation}
For fixed $k$ there are at most $(k!)^{|\mathcal A|}$ discrete choices (a crude bound before simultaneous conjugacy reduction), and the remaining feasible set is compact. The optimal stationary error at any fixed width is attained.
\end{theorem}
\begin{proof}
Necessity follows from Theorem~\ref{thm:purification55}, padding by unused labels if its output has fewer than $k$ states, and then taking the closure of word products. Sufficiency is the restriction of \eqref{eq:permutationfeasibility55} to words. The probabilities and legal decoders lie in a finite product of compact sets. Each displayed inequality defines a closed subset of that product. The worst error is continuous in these parameters, uniformly in the compact group variable, so its minimum exists. There are finitely many permutation tuples.
\end{proof}

\begin{theorem}[Finite observable quotient criterion]\label{thm:finitequotient55}
For arbitrary compact $H$ and the stated continuous interface,
\[
 M_\varepsilon<\infty
 \quad\Longleftrightarrow\quad
 r(U)\le\varepsilon\text{ for some open normal subgroup }U\le H.
\]
If such $U$ exists, $|\mathcal S|[H:U]$ deterministic-initialized permutation labels suffice. Conversely, a stationary $k$-label realization supplies such a $U$ with $[H:U]\le k!$.
\end{theorem}
\begin{proof}
Purify a stationary realization and use its joint group $\mathcal L_\sigma$. Put
\[
                  U=\{h\in H:(h,I)\in\mathcal L_\sigma\}.
\]
This is a closed normal subgroup of $H$: conjugate $(h,I)$ by an element above any $g\in H$, using surjectivity of the first projection. For a fixed permutation $\Pi$, a nonempty fiber in $\mathcal L_\sigma$ is a coset $hU$. These at most $k!$ fibers cover $H$, so $U$ has finite index. A closed subgroup of finite index in a compact group is open, since its complement is a finite union of closed cosets.

For each $s,j$ and each fiber, the single legal vector $\beta_s\Pi d_j$ is within $\varepsilon$ of every $F_{s,j}(h^{-1})$ in that fiber. Explicitly, $(hU)^{-1}=Uh^{-1}=h^{-1}U$ by normality, so inversion preserves the ordinary $U$-coset form. Hence $r(U)\le\varepsilon$. Conversely, store $(s,gU)$, update the coset by left multiplication by $u_a$, and decode a minimizing center of $F_{s,j}(gU)$. Compactness supplies the centers. This gives the asserted finite realization.
\end{proof}
For finitely many connected components, every open subgroup contains $K=H^\circ$, and $K$ itself is open normal. Therefore $M_\varepsilon<\infty$ holds exactly when
\[
 \varepsilon\ge\max_{s,j,c\in H/K}\rad_{Y_j}(F_{s,j}(c)),
\]
\emph{without} a return assumption. Executable returns enter the stronger clocked converse, not this stationary statement. The index bound $k!$ is an existence bound on a deterministic quotient, not a claim that the stationary minimum equals its index.

\begin{proposition}[Random initialization strictly saves states]\label{prop:C655}
Let $H=\Z/6\Z$, with an identity letter and a generator, one seed, and binary success probability
\[
 f(n)=\tfrac12+\tfrac14(-1)^n+\tfrac14\cos(2\pi n/3).
\]
At zero error the least deterministic component-quotient width is six, whereas the unrestricted stationary stochastic minimum is five.
\end{proposition}
\begin{proof}
The values at $n=0,\ldots,5$ are
\[
                  (1,\tfrac18,\tfrac58,\tfrac12,\tfrac58,\tfrac18).
\]
This sequence has least period six. A deterministic quotient of the cyclic component action preserving these outputs must therefore have six states. For a five-state realization take the disjoint union of a two-cycle and a three-cycle, both updating the index by $i\mapsto i+1$ modulo their respective lengths, initialize with mass $1/2$ at the zeroth vertex of each, and use decoder rows $(1,0)$ on the two-cycle and $(1,1/4,1/4)$ on the three-cycle. The identity command fixes all five labels. Direct averaging gives exactly $f(n)$, with legal binary probabilities.

If a stationary realization on at most four labels existed, purification would give one whose generator is a permutation of at most four letters. Its order is one of $1,2,3,4$, as follows by listing the possible cycle lengths. Every resulting numerical sequence has a period dividing that order, contradicting the least period six. Thus five is the unrestricted minimum.
\end{proof}
